\paragraph{Bilingual interval learning and question assignment.}
For pictured page positions $i\le j<i+48$, a bidirectional reader assigned the score $u_i+v_j$. Its loss was the log sum of exponentiated scores over all valid intervals minus the corresponding log sum over all human-accepted occurrences. The 11,438 parallel training pairs retained complete within-page exchange cycles. Both runs updated the last four contextual layers and endpoint projection of the 86.0-million-parameter visually adapted Chinese QA encoder, exposing each Chinese and English request once. The exchanged control assigned English questions to the other facts on the same page. The optical encoder and the first eight contextual layers were frozen. The decoder returned the single highest-scoring interval. Several locations can have the same physical spelling, so this decision can differ from selecting the spelling with greatest total mass.

Three classical works and two modern essays supplied 23 further passages. Each had two disjoint extractive answers and two question languages. Source and question images used the same regular typeface as the fit. Scoring required an accepted complete bitmap sequence. For the writing test, the released causal writer received either the original page or a predicted excerpt, together with the original question. Writer scoring also required an autonomous stop.

A separate 0.85-million-parameter model encoded pictured questions with a shared optical projection, one-dimensional convolutions, two attention layers and learned pooling. Eight passes used symmetric contrastive likelihood with temperature 0.07. Each batch included complete groups of different facts from the same training page. If a source or target question duplicated another pictured question in the batch, its corresponding alignment also counted toward the positive mass. Independent matched and exchanged runs used identical initial weights. If $s_{ij}$ is the cosine similarity between English question $i$ and Chinese question $j$, a strict assignment requires $s_{11}+s_{22}>s_{12}+s_{21}$. Additive row or column terms cancel exactly. Lean checks this identity and a perturbation bound for integer score units. The empirical scores use floating-point cosines.
